\documentclass[11pt,a4paper]{article}

\usepackage[margin=1in]{geometry}
\usepackage{hyperref}
\usepackage{amsmath,amssymb}
\usepackage{xcolor}
\usepackage{parskip}

\newcommand{\rev}[1]{\medskip\noindent\textit{#1}\par\medskip}
\newcommand{\reply}{\noindent\textbf{Response.}\ }

\title{Response to Reviewers\\
\large High-Level Structure of Dependency Networks}

\author{Alexei Vazquez}
\date{}

\begin{document}

\maketitle

I thank the three reviewers for their careful and constructive reading. The
comments identified two genuine weaknesses --- insufficient statistical support
for the claimed relations, and an incomplete engagement with the neighbouring
literatures --- and several presentational problems. I have addressed all of
them, and in doing so I have had to weaken two of the claims made in the
original manuscript. Specifically:

\begin{itemize}
\item The number of projects analysed is 1198, not a handful. This was not
stated in the original manuscript, which is why Reviewer 1 reasonably assumed
the sample was small. A descriptive statistics table has been added.
\item Formal model selection now supports the logarithmic form for the DUS and
the sub-linear power law for the DPS, but the DUS relation explains only 27\%
of the variance and does not extrapolate to the public PSPLIB benchmark. I have
therefore stopped calling it a scaling law and describe it as slow logarithmic
growth of the central trend.
\item The small-world question is now settled by the analysis Reviewer 1 asked
for. Clustering is given a DAG-appropriate definition, two degree-preserving
null models are built, and the result is that the short distances follow from
the degree distributions of the DUS network rather than from any special
arrangement of the shortcuts. The distances are short; the Watts--Strogatz
criterion is simply not the informative test for these networks.
\item The finer construction has been renamed from DFS to DPS (Dependency
Predecessor Structure) to remove the clash with depth-first search.
\item All key results are now replicated on public data (PSPLIB and two
anonymised real schedules shipped with the code), and every number and figure
in the paper is produced by scripts included in the public repository.
\end{itemize}

The paper has been restructured accordingly: the schematic that was Figure 4 is
now Figure 1, a descriptive statistics table and two new figures have been
added, and the manuscript now has six figures and one table.

\subsection*{A correction found while preparing this revision}

While rebuilding the analysis so that it could be released publicly, I
recomputed every quantity directly from the raw schedules and compared the
result against the summary table used for the original submission. The number of
activities, the number of dependencies, the number of DUS levels and the
distances agreed exactly for all 1198 projects. The number of DPS classes did
not.

The cause was a defect in the original code. The DPS label packs one bit per DUS
level, and it was accumulated with \texttt{math.pow(2,\,j)}, which returns a
float. A float64 mantissa holds 53 bits, so once a schedule has more than 53 DUS
levels the low-order bits are lost, distinct predecessor patterns collapse onto
the same value, and counting distinct labels under-estimates the number of DPS
classes. The error is exactly zero at or below 53 levels and grows with depth:
across the 1198 projects the under-count had a median factor of $1.010$, rising
to a median of $1.05$ between 100 and 200 levels and $1.17$ above 200 levels,
with a worst case of $1.51$.

Because the under-count grows with project size, it biased the fitted exponent
downward. Every DPS result in this revision has been recomputed with exact
integer arithmetic. The corrected exponent is $0.709$ $[0.692, 0.724]$ rather
than $0.694$ $[0.679, 0.709]$, the prefactor is $1.42$ rather than $1.56$, and
the DPS column of Table 1 shifts slightly. The conclusion is
unchanged and, if anything, strengthened: the growth is a power law, decisively
preferred over the alternatives, with an exponent well below $1$. The DUS
results, the distance results and Figures 3 and 4 are untouched, since the
defect affected only the DPS labelling.

I report this rather than silently substituting the corrected numbers. The fix
is committed to the public repository with the diagnosis in the commit message,
and the released implementation now builds the label with exact integer bit
shifts. It was Reviewer 1's request for public replication that prompted the
rewrite which exposed the problem.

\section*{Reviewer 1}

\rev{The manuscript proposes a deterministic two-step coarse-graining method for
dependency DAGs\ldots{} the current version of this paper shows some weaknesses in
terms of its originality, statistical rigor, and the validation of the two
arguments presented.}

\reply I agree with this assessment of the original version. Each of the four
points is addressed below.

\subsection*{1. Missing comparison with DSM, WBS/CPM and DAG condensation}

\rev{From a project management perspective, the Design Structure Matrix (DSM)
literature has a long history regarding the clustering and ordering of dependency
matrices, and the concepts of WBS and CPM are conceptually very similar. Yet,
this paper does not cite or compare these studies. From a computer science
perspective, DAG condensation and hierarchical graph techniques are similarly
relevant.}

\reply Accepted. The Introduction now contains a paragraph that positions the
work against both literatures and states explicitly what is and is not new.

On the project management side, I now cite Steward's original DSM paper,
Browning's review and the Eppinger--Browning monograph, and I draw the
distinction that matters: DSM sequencing and the layering used here both look
for an ordering compatible with the dependencies, but DSM methods are typically
applied to matrices that contain cycles, where the object is to find and break
feedback loops, and DSM clustering is usually the outcome of an optimisation
with tunable objectives and no unique answer. The construction here applies to
acyclic networks, involves no optimisation and no free parameters, and is aimed
at compression rather than resequencing. CPM is now cited at the point where the
schedule is introduced, and the relationship is made explicit in the DUS section:
the DUS index of the final activity is the number of activities on the longest
dependency chain, which is the critical path when all activities have equal
duration. The relation to the WBS is now a structural part of the paper rather
than a passing remark --- Figure 1A shows a WBS, Figure 1B the dependency network
of its leaves, and the caption points out a dependency that crosses a WBS branch
boundary, which is the reason the two views do not coincide.

On the computer science side I now cite and distinguish three things that might
otherwise be confused with this construction. Condensation by strongly connected
components (Tarjan) is the standard way to turn a digraph into a DAG; applied to
a schedule, which is already acyclic, it does nothing at all. Transitive
reduction (Aho, Garey and Ullman) removes redundant edges but preserves every
node and therefore achieves no node-level compression. Layered graph drawing
(Sugiyama; Healy and Nikolov) does assign nodes to layers by essentially
Eq.~(1), which the original manuscript already acknowledged, but treats the
layering as a step towards a picture rather than as an object of study.

I want to be straightforward about the resulting claim of novelty, which is
narrower than the original manuscript implied. The DUS index is the standard
longest-path layering and I do not claim it as new. What I claim is: its use as
a coarse-graining with measured compression, the empirical relations of Figures
3 and 5 with the model selection now supporting them, the DPS refinement, and
the observation that the resulting networks are shallow. This is now stated in
those terms at the end of the DUS section.

\subsection*{2. Statistical support for the scaling relations}

\rev{(i) The logarithmic fit in Figure 2 appears to be a single-parameter fit;
the sample size, a goodness-of-fit measure, and a confidence interval on p should
be reported\ldots{} a model-selection argument (e.g., AIC/BIC or a likelihood-based
comparison) is needed to justify the logarithmic form.}

\reply Accepted in full, and this changed the results. The analysis has been
redone as follows. Four functional forms (one- and two-parameter logarithms, a
power law, a linear relation) were each fitted by maximum likelihood under two
error structures (additive and multiplicative), giving eight fits. The
multiplicative log-likelihoods include the Jacobian term $-\sum_i \ln y_i$ so
that all eight are likelihoods of the same observed quantity and their AIC
values are directly comparable across both functional form and error structure.
This is described in a new Statistical methods subsection.

The outcome supports the logarithmic form, but with a caveat the original
manuscript did not report:

\begin{itemize}
\item The preferred model is the two-parameter logarithm with multiplicative
errors, $m = 14.3\ln x - 33.9$. The reviewer was right that the published fit
was a one-parameter logarithm; that form is now rejected by
$\Delta\mathrm{AIC} = 45$.
\item It is preferred over the power law by $\Delta\mathrm{AIC} = 22$ and over
the linear relation by $\Delta\mathrm{AIC} = 147$. The best-fitting power law
has exponent $0.17$, which as the reviewer anticipated is small enough that the
two forms are hard to separate visually; the likelihood comparison is what
separates them.
\item $n = 1198$. The 95\% bootstrap confidence interval on the logarithmic
coefficient is $[13.4, 15.2]$. Residuals show no detectable heteroscedasticity
(Breusch--Pagan $p = 0.07$).
\item Goodness of fit is modest: $R^2 = 0.27$ on the logarithm. This is now
stated in the text, and Figure 3 shows binned medians with interquartile ranges
so that the reader can see the scatter directly.
\end{itemize}

Given $R^2 = 0.27$, and given the extrapolation failure described under point 4
below, I have removed the term ``scaling law'' for this relation and describe it
instead as slow logarithmic growth of the central trend. The compression claim
survives and is now quantified concretely: a tenfold increase in dependencies
adds about 33 DUS levels regardless of starting size.

\rev{(ii) The power-law fit in Figure 5 should specify the fitting method; a naive
log-log OLS is biased, so a Clauset--Shalizi--Newman-style procedure (Clauset et
al. 2009), or at least residual diagnostics and a confidence interval on the
exponent, would be appropriate.}

\reply The fitting method is now specified in full, and residual diagnostics and
a confidence interval are reported. The preferred model is a power law under
multiplicative errors, $1.42\,x^{0.71}$, ahead of the linear relation by
$\Delta\mathrm{AIC} = 269$ and the two-parameter logarithm by
$\Delta\mathrm{AIC} = 1310$. The exponent is $0.709$ with a 95\% bootstrap
confidence interval of $[0.692, 0.724]$, comfortably excluding $1$, and
$R^2 = 0.81$ on the logarithm. The residuals of this fit are heteroscedastic,
which is now stated explicitly; this is precisely why the interval quoted is a
nonparametric case-resampling bootstrap rather than a nominal standard error.

On the Clauset--Shalizi--Newman procedure I would respectfully disagree, and
have added a short paragraph to Methods explaining why. That procedure fits and
tests a heavy-tailed \emph{probability distribution} of a single random
variable, using maximum likelihood for the tail exponent above an estimated
$x_{\min}$ together with a Kolmogorov--Smirnov goodness-of-fit test. The
relation in Figure 5 is not a distribution: it is a regression of one measured
quantity (number of DPS classes) on another (number of dependencies) across
projects. The estimator, the $x_{\min}$ selection step and the KS test all
presuppose i.i.d. draws from a distribution and have no direct analogue here.
The reviewer's underlying concern --- that naive log-log OLS is not a neutral
choice --- is entirely valid, and I have addressed it in the way that is
appropriate for a regression: by making the error model an explicit part of the
model comparison rather than an unstated assumption, and by using a bootstrap
that does not assume homoscedasticity. I would of course reconsider if the
reviewer thinks a distributional formulation is called for.

\rev{(iii) The claim that all three companies follow the same relation should be
supported by a test of parameter commonality rather than visual impression. When
the sample size is small (only a few projects per company), the use of strong
terms such as ``scaling law'' should be used with appropriate moderation.}

\reply Accepted, and this is now tested formally. I should first correct a
factual premise, and the fault for it is mine: the sample is not small. It
comprises 1198 projects --- 82, 249 and 867 from the three companies. The
original manuscript said only that ``each company contributed several
projects'', which was uninformative and invited exactly this reading. The new
Table 1 gives the full breakdown.

Commonality is tested with nested likelihood ratio tests comparing three models:
fully pooled parameters; a shared logarithmic coefficient or power-law exponent
with company-specific offsets or prefactors; and fully company-specific
parameters. This separates the question of whether the companies share the
\emph{form} of the relation from whether they share its nuisance parameters, and
the two answers turn out to be different:

\begin{itemize}
\item DUS: a shared logarithmic coefficient is not rejected
($\mathrm{LR} = 2.4$, $\mathrm{df} = 2$, $p = 0.30$; shared value $15.1$), and
AIC prefers the shared-coefficient model. The offsets do differ
($p = 5 \times 10^{-9}$).
\item DPS: a shared exponent is not rejected ($\mathrm{LR} = 1.1$,
$\mathrm{df} = 2$, $p = 0.59$; shared value $0.72$), while the prefactors differ
($p = 4 \times 10^{-6}$).
\end{itemize}

So the original claim that ``the same relation describes all three companies''
is supported in the sense that matters --- the rate of growth is common --- but
was too strong as stated, because the intercepts and prefactors are
significantly company-specific. The revised text says exactly this and
interprets the difference as company-specific scheduling practice shifting the
absolute level without changing the rate.

On terminology I have followed the reviewer's advice throughout. The phrase
``scaling law'' no longer appears anywhere in the manuscript as a claim. The DUS
relation is described as slow logarithmic growth of the central trend, for the
reasons above. The DPS relation is described as a sub-linear power law, which is
a statement about functional form rather than a claim of universality, and it is
the better supported of the two: a power law is preferred by a very large
margin, the exponent is tightly determined, $R^2 = 0.81$, and the relation
extrapolates successfully to independent public data.

\subsection*{3. The small-world claim}

\rev{The evidence supporting the small-world claim is weak\ldots{} this argument
simply shows that the average distance is less than the expected value for a
pure chain. Since any shortcut reduces the distance compared to a chain, this is
almost a trivial statement. The Watts--Strogatz definition requires not only
short path lengths but also high clustering compared to degree-preserving random
graphs\ldots{} the authors would need to clearly specify how clustering is defined
in a DAG.}

\reply Accepted on both counts, and the analysis the reviewer asks for has now
been carried out. Falling below the chain expectation is indeed nearly
tautological, and the original manuscript did not make the clustering
comparison. The reviewer is also right that clustering needs defining first: the
standard coefficient counts closed triangles, and a {\em cyclic} triangle cannot
exist in a DAG.

That last point has a clean resolution, which the revision adopts. A cyclic
triangle is impossible in a DAG, but a {\em transitive} one is not: $i \to j$,
$j \to k$ and the shortcut $i \to k$ are mutually consistent with acyclicity. We
therefore use the transitive clustering coefficient $C$, the fraction of
directed two-paths $i \to j \to k$ closed by an edge $i \to k$. It is the
admissible member of the directed-triangle family, it is precisely the quantity
that a shortcut creates, and it is zero for a pure chain. This is now stated in
the Results and defined in a new Methods subsection.

Two degree-preserving null models were then built, both by double-edge swaps.
The first holds the DUS spine fixed, since Eq.~(1) forces it and randomisation
should not be free to destroy it, and rewires the skip dependencies while
preserving every level's skip in- and out-degree. The second preserves the in-
and out-degree of every {\em activity}, keeps the network acyclic by requiring
every edge to point forward in a fixed topological order, and re-derives the DUS
decomposition from scratch. Mixing was checked in both cases.

The result, over 198 projects, is that the observed clustering and the observed
distances are almost exactly those of the first null: $C/C_{\mathrm{null}} =
1.04$ (IQR $[1.01, 1.10]$) and $L/L_{\mathrm{null}} = 1.03$ (IQR
$[1.00, 1.08]$). This is a new Figure 7 and a new Results subsection.

I want to be careful about what this does and does not show, because it is easy
to misread it as a negative result. It does {\em not} cast doubt on the
networks having small average distances --- that is a measured property of the
data, and no null model can remove it: the mean distance to the end is $2.5$
levels where a chain of the same depth would give up to $285$. What it shows is
{\em why} the distances are small. Given how many long-range dependencies each
level has, placing them anywhere at all reproduces the observed distances, so
the shallowness follows from the degree distributions of the DUS network rather
than from any special arrangement of the shortcuts. That is a mechanistic
explanation, and a positive one. It also connects to the generative side of the
problem: the duplication-split model of activity networks (Vazquez et al.,
Sci.\ Rep.\ 13, 509, 2023) produces degree distributions of exactly this kind,
and the present result establishes that such distributions are sufficient for
the effect.

What the analysis does bear on is the Watts--Strogatz {\em criterion}
specifically, which asks for clustering substantially in excess of a
degree-matched null. That excess is absent, so the criterion is not the
informative test for these networks, and the manuscript states the property
directly rather than invoking the label. This is a statement about which
diagnostic is appropriate, not about whether the distances are short.

The second null adds a complementary finding. Preserving only the activity
degree sequence and re-deriving the decomposition gives DUS networks roughly
half as deep as the real ones (median $m/m_{\mathrm{null}} \approx 2$, a
conservative figure, since this randomisation mixes slowly and the gap widens
as it proceeds). So the degree distributions determine how short the paths
through the DUS network are, but not how deep it is; real schedules retain
genuine sequential structure that the degree sequence alone does not capture.
The Discussion now poses the generative origin of that depth as the open
question, in place of the definition of clustering, which is settled above.

The distance evidence itself has also been strengthened beyond the chain
comparison, because the interesting fact turned out not to be the gap but the
flatness:

\begin{itemize}
\item The mean distance to the final DUS level has median $2.5$ with
interquartile range $[1.9, 3.6]$, essentially independent of the number of
levels $m$, which ranges from $14$ to $571$. Binned medians are $3.4$, $2.3$ and
$2.7$ for small, medium and large projects.
\item The rank correlation between distance and $m$ is $\rho = 0.10$
($p = 8\times10^{-4}$, $n = 1198$): nominally significant at this sample size
but negligible in magnitude. For completeness the manuscript also reports that
the same model selection procedure formally prefers a weakly increasing linear
relation over a strictly constant one ($\Delta\mathrm{AIC} = 20$, slope $0.0041$
per level, 95\% CI $[0.0018, 0.0062]$). The implied growth is negligible: across
the whole range the fitted distance rises from $2.5$ to $4.7$, a factor of
$1.9$, while $m$ grows by a factor of $41$. I preferred to report this rather
than claim exact constancy.
\item Figure 4 is now on logarithmic axes, which makes the contrast between the
flat data and the rising chain expectation directly visible; the original linear
axes obscured it.
\end{itemize}

The word ``small-world'' now appears only where the shortcut mechanism is
attributed to Watts and Strogatz, and where the Watts--Strogatz criterion is
discussed and set aside for the reason given above, never as an unexamined
label.

\subsection*{4. Proprietary data, descriptive statistics and public replication}

\rev{It would be advisable for the paper to include at least descriptive
statistics---namely, the number of projects per firm, the range of activity and
dependency counts, and basic summary information on the distribution.
Furthermore, reproducing key scalability results using publicly available
scheduling datasets (e.g., PSPLIB) would significantly enhance the replicability
of the research arguments.}

\reply Both accepted and done.

Table 1 now reports, per company and overall, the number of projects and the
median and full range of the number of activities, dependencies, DUS levels and
DPS classes. The dataset spans 37 to 104{,}981 activities and 36 to 150{,}692
dependencies. The Methods section also now records the one schedule that was
excluded from the original analysis and why (its first DUS level connects to
every other level, so its DUS network is a star rather than a network with a
spine); this exclusion was applied but not reported in the original manuscript.

For replication I have analysed all 2040 single-mode PSPLIB instances (j30, j60,
j90, j120), plus the two anonymised real schedules that are distributed with the
code repository and which extend the publicly reproducible range to $1.3\times
10^{5}$ dependencies. This is a new figure (Figure 6) and a new Results
subsection. The result is mixed, and I report it as such:

\begin{itemize}
\item PSPLIB instances have 32--122 activities and 48--257 dependencies, come in
only four activity counts, and span barely half a decade in size. Model selection
within PSPLIB alone favours a power law for the DUS and a linear relation for the
DPS, which are not the forms selected on the construction data; over so narrow a
range I do not treat that as evidence about functional form, and the manuscript
reports it as such rather than presenting a within-PSPLIB fit as a confirmation.
\item What PSPLIB can test is transfer. Its dependency range overlaps only the
extreme lower tail of ours --- just $2.4\%$ of the construction projects are that
small --- so it probes the relations where they were barely constrained. The DPS
power law, fitted on construction data and evaluated two orders of magnitude
below the median project on which it was estimated, predicts the PSPLIB values to
a median observed/predicted ratio of $0.77$. The power law fitted within PSPLIB
alone gives exponent $0.86$, also sub-linear. This is a genuine out-of-sample
success.
\item The DUS logarithm does not extrapolate: it over-predicts PSPLIB by more
than a factor of two (median ratio $0.42$). The paper now states that the
logarithmic relation describes large construction schedules rather than being
universal, which is the second reason for dropping the term ``scaling law'' for
it.
\item The qualitative shallowness result does replicate: median distance to the
final level $2.1$, below the chain expectation in $100\%$ of instances.
\end{itemize}

Finally, the code repository now contains the scripts that perform every
statistical test, download and process PSPLIB, and generate every figure and the
table in the paper, so that all public-data results can be reproduced end to end.

\section*{Reviewer 2}

\rev{Both methods reveal a small-world structure in the underlying dependency
network\ldots{} I would like to highlight two main limitations.}

\reply I thank the reviewer for an accurate summary. Note that the small-world
claim has been withdrawn following Reviewer 1's point 3; see that response.

\subsection*{1. Structure and placement of illustrative material}

\rev{The structure of the paper is unconventional, and key illustrative material
is placed too late. For instance, the illustration of the two metrics appears in
Figure 4, whereas it would be far more effective as Figure 1, orienting the
reader from the outset.}

\reply Accepted. The schematic is now Figure 1 and appears in the Introduction,
before either construction is defined. It has also been substantially redrawn
and expanded to three panels so that it can carry more of the exposition:
panel A is a small Work Breakdown Structure, panel B the dependency network of
its leaf activities with the DUS index shown as vertical bands and the DPS class
shown by colour, and panel C the two coarse-grained networks. The example was
chosen so that the DPS genuinely refines the DUS --- two activities share a DUS
level and a DPS class while a third shares their level but forms its own class
--- which the original schematic did not illustrate. Panels A and B also make
visible the point that dependencies cross WBS branches.

The figure order is now Figure 1 schematic, Figure 2 real DUS profiles, Figure 3
DUS scaling, Figure 4 distances, Figure 5 DPS scaling, Figure 6 public
replication, plus Table 1.

\rev{I would introduce DAGs more extensively, and move the first paragraph of the
Results section into the Introduction.}

\reply Both done. The first paragraph of Results has been moved into the
Introduction and expanded into a fuller treatment of DAGs: the definition, the
acyclicity requirement and why it is forced, topological order, the notion of
incomparability that the DUS relies on, and a note that DAGs also represent
causal models, ancestries and feedforward networks, which is what makes the
constructions transferable. The Results section now opens directly with the DUS
definition.

\subsection*{2. Missing literature: symmetry, compression, hyperbolic embeddings}

\rev{The authors overlook a substantial body of related work, including directed
acyclic graphs (DAGs), network symmetry, compression symmetry, and hyperbolic
embeddings. As examples, I would point to Krioukov et al., ``Hyperbolic geometry
of complex networks'', and S\'anchez-Garc\'ia, ``Exploiting symmetry in network
analysis''.}

\reply Accepted, and both suggested references are now cited, along with
related work the reviewer's comment led me to. The new related-work paragraph in
the Introduction places this paper among approaches to meso-scale structure that
do not rely on statistical inference: exact symmetry (MacArthur, S\'anchez-Garc\'ia
and Anderson; S\'anchez-Garc\'ia), graph fibrations (Boldi and Vigna; Morone et
al., already cited in the original), latent hyperbolic geometry (Krioukov et
al.), and explicit hierarchy measures for directed networks (Mones et al.;
Corominas-Murtra et al.).

The connection to the symmetry literature is more than bibliographic and I have
made it explicit in the DPS section. Two activities in the same DPS class have
predecessor sets occupying the same DUS levels, which is a local symmetry of the
sort that fibrations formalise; the DPS can be read as a coarse, level-indexed
relaxation of a fibration, which is what makes it computable by a single pass
and labellable by one integer. S\'anchez-Garc\'ia's framing of symmetry as a
route to network reduction is exactly the use made of it here.

\subsection*{Minor comments}

\rev{The construction of the DUS and DFS networks needs clearer explanation;
there is currently some confusion between the network itself and the associated
metric/category.}

\reply Accepted; this was a real ambiguity in the original text, which used
``DUS'' for the index, for a single level, and for the partition. The paper now
names four distinct objects and uses them consistently: the \emph{DUS index}
$u_i$, an integer attached to an activity; a \emph{DUS level}, the set of
activities sharing one index value; the \emph{Dependency Uprise Structure}, the
partition into levels; and the \emph{DUS network}, whose nodes are levels. The
same four-way distinction is drawn for the DPS. Phrases such as ``the number of
DUSs'' have been replaced throughout by ``the number of DUS levels'' or ``the
number of DPS classes''.

\rev{``Equation (1)'' is referred to elsewhere in the text simply as ``(1),''
which is ambiguous on first read.}

\reply Fixed. All equation references now read ``Eq.~(1)'' or ``Eq.~(2)''.

\rev{Figure 1 is very difficult to read due to the complexity (number of nodes
and edges) of the network, which obscures a clear understanding of the method's
benefits.}

\reply Accepted, and the figure has been redrawn (it is now Figure 2). The
original used a semicircular arc layout in which the level sizes, the spine and
the skip dependencies were superimposed in one crowded picture. The new version
separates them: the DUS spine runs along a linear horizontal axis, bars below
give the number of activities per level coloured by progress, and arcs above
show dependencies between non-adjacent levels, drawn as semicircles whose radius
is the span so that long-range links are visually distinct from short ones. Only
the strongest 5\% of skip dependencies are drawn, with the total stated on the
figure so that nothing is hidden silently. Panel titles now give the activity,
dependency and level counts. Both schedules shown are included in the public
code repository, so the figure can be regenerated by any reader.

\section*{Reviewer 3}

\rev{Overall I like the paper, but I have a few comments.}

\reply I thank the reviewer for the encouragement and for a set of comments that
improved the exposition considerably.

\subsection*{1. An example of a WBS}

\rev{I am not in project management, so in the introduction it would be useful to
see an example (even a simplified one) of a WBS, to get a better sense of the
objects of study.}

\reply Done. Figure 1A is a small WBS for a house, decomposed into substructure,
superstructure and services and then into seven schedulable activities, and the
Introduction now walks through it. The text notes that real construction WBSs
have exactly this shape but reach tens of thousands of leaves. Figure 1B then
shows the dependency network over those same seven activities, so the reader can
see both objects and the fact that they do not coincide.

\subsection*{2. A diagram for the DUS, and the relation to breadth-first search}

\rev{The definition of the DUS is clear enough, but again, it seems like a diagram
would be helpful here. Perhaps move Figure 4 to this section. How does the DUS
index differ from a breadth first search starting from the root `project
completed' node?}

\reply The diagram has been moved forward as requested and now appears as
Figure 1, ahead of both definitions; it is referenced from the DUS section.

The breadth-first search question was a good one and the answer is now a
paragraph in the DUS section. The two differ because breadth-first search
computes \emph{shortest} path lengths while Eq.~(1) computes \emph{longest}
ones. They disagree whenever an activity has predecessors at different depths.
The paper gives the concrete case from Figure 1B: a breadth-first search from
the start would place Utilities at depth 1, because Excavate feeds it directly,
whereas Eq.~(1) places it at $u = 2$ because it must also wait for Foundations.
Only the longest-path rule guarantees that every one of an activity's
predecessors lies strictly below it, which is the property the whole
construction depends on --- it is what makes the levels ordered, what makes
activities within a level pairwise incomparable, and what guarantees the DUS
spine exists. A backwards breadth-first search from the completion node is also
addressed: it would group activities by remaining depth to the end rather than
by accumulated depth from the start, which is a different and in general
incompatible grouping.

\subsection*{3. The name ``uprise''}

\rev{On naming, what does `uprise' mean in this context? The usual use of the word
is as a verb `uprising' meaning a rebellion which I don't think is the
intention! Why not just `dependency structure'?}

\reply The reviewer is right that the word is unusual, and the intended sense
was never stated. It is now: the index counts how far an activity has risen
above the start of the project, in dependency steps rather than in calendar
time, and a sentence in the DUS section says so.

I have kept the name, for two reasons I hope the reviewer will find reasonable.
``Dependency structure'' on its own would be ambiguous, since it is a natural
description of the dependency network itself, and it also collides with Design
Structure Matrix terminology, which Reviewer 1 has asked me to engage with
directly. And the name and its acronym are already in use in the public code
repository. If the editor or the reviewer feels the term remains an obstacle, I
am willing to rename it --- ``Dependency Level Structure'' would be the natural
alternative --- but I did not want to change it unilaterally given the acronym
is now load-bearing across the paper and the software.

\subsection*{4. Redundant panels in Figure 1}

\rev{Figure 1 is nice, but why include A and B when the same information is in C
and D?}

\reply Accepted; panels A and B have been removed. The figure (now Figure 2)
shows only the two annotated profiles, which as the reviewer observed contain
everything the bare networks did. The freed space has been used to draw the
remaining panels larger and to add the readability improvements requested by
Reviewer 2.

\subsection*{5. The flat trend in Figure 3}

\rev{I realise the point of Figure 3 is to show real DUSs fall below the line, but
the data appears mostly constant indicating the same average distance to the end
regardless of project size? Perhaps this is obvious in project management but
seems worth commenting on.}

\reply This was the single most useful observation in the reviews, and the
reviewer is right. The data are nearly constant, this is not obvious in
project management, and it is a stronger and more interesting statement than the
one the original manuscript made. It is now the main claim of that section rather
than an unremarked feature of the plot.

Quantitatively: median distance $2.5$, interquartile range $[1.9, 3.6]$, with
binned medians of $3.4$, $2.3$ and $2.7$ for projects with fewer than 50, 50--150
and more than 150 DUS levels, while $m$ itself ranges from $14$ to $571$. The
rank correlation with $m$ is $\rho = 0.10$, nominally significant at $n = 1198$
but negligible. Figure 4 is now plotted on logarithmic axes with binned medians
overlaid, which makes the flatness against the rising chain expectation obvious;
on the original linear axes it was not. The text draws the practical
consequence: a schedule with 500 DUS levels is not 500 steps deep in any
operational sense, because long-range dependencies carry information to the end
in two or three steps. This also connects to Reviewer 1's point 3, where the
strong small-world claim has been replaced by exactly this statement.

\subsection*{6. The DFS name and DPS visualisation}

\rev{Regarding the name, I strongly suggest you change it since DFS=Depth First
Search and this will probably cause a lot of confusion. Does the DFS admit any
nice visualisations like Figure 1?}

\reply Accepted without reservation --- the clash is real and would have caused
confusion, particularly since the paper discusses graph traversal explicitly in
answering the reviewer's point 2. The construction has been renamed throughout
to \textbf{Dependency Predecessor Structure (DPS)}. The new name is also more
descriptive: the grouping is defined precisely by which DUS levels an activity's
predecessors occupy. The connection to graph fibrations is retained and made
more explicit in the text, but it no longer drives the acronym.

On visualisation, the DPS is now shown in Figure 1B, where DPS class is encoded
by node colour, and in Figure 1C, where the DPS network is drawn alongside the
DUS network so that the two coarse-grainings can be compared directly on the
same example. I did not add a DPS analogue of the real-schedule profile in
Figure 2. The reason is that a DPS network of a real project has of order
$10^3$ nodes with no canonical linear ordering --- the DPS index is a
categorical label, not a coordinate --- so the arc-diagram layout that makes the
DUS profile legible has nothing to lay out against. Finding a good visual
representation for the DPS network is noted as future work.

\subsection*{7. Insight for project management}

\rev{I assume the purpose of these constructions is to provide some insight into
project management schedules. The author should discuss this.}

\reply Accepted; the Discussion now has a paragraph devoted to this, giving
three concrete uses. Reporting and tracking: the DUS profile compresses a
schedule of tens of thousands of activities into a picture with of order a
hundred nodes on which progress, cost or risk can be overlaid, and in which a
stalled level or an unusual concentration of work is immediately visible.
Diagnosis: the long arcs in Figure 2 are dependencies that skip many levels and
therefore propagate delay furthest, which makes them natural candidates for
scrutiny --- such a link is usually either a real structural constraint or an
artefact of how the schedule was built. Input to risk models: a coarse-grained
network with a hundred nodes is tractable for Monte Carlo simulation where one
with a hundred thousand is not. The paragraph also notes why determinism matters
in this setting: two planners applying the construction to the same schedule get
the same answer, which is not true of clustering methods that depend on an
objective function or a random seed. A further paragraph now states three
limitations explicitly, including that the DUS is one antichain decomposition
among many and is not claimed to be optimal by any criterion.

\end{document}
